\subsection{Third Construction: The Perpendicular Line}

\begin{derivationbox}[Third construction]
Through $O$, construct the unique line perpendicular to the first line. The two
constructed lines now intersect at a right angle.
\end{derivationbox}

\begin{definitionbox}
We call the first line the $x$-axis, the perpendicular line the $y$-axis, and
their intersection $O$ the origin.
\end{definitionbox}

\begin{center}
\begin{tikzpicture}[scale=0.9]
    \draw[axis] (-4.2,0) -- (4.4,0) node[right] {$x$};
    \draw[axis] (0,-3.1) -- (0,3.3) node[above] {$y$};
    \fill (0,0) circle (2pt);
    \node[below left] at (0,0) {$O$};
    \draw (0.26,0) -- (0.26,0.26) -- (0,0.26);
\end{tikzpicture}
\end{center}

\begin{corebox}[What has now been created]
Only after the perpendicular line has been constructed do we have the geometric
framework of two coordinate axes with a common origin.
\end{corebox}

\begin{remarkbox}
The axes are perpendicular because this was required in their Euclidean
construction. No coordinate formula and no numerical angle measure was used to
create that right angle.
\end{remarkbox}
